\section{The packaging instrument}\label{sec:instrument}

This section fixes the object the rest of the paper studies: set theory read as a Six Birds
\emph{theory package}~\citep{SBTFoundationsI}. We keep the framework's apparatus to the minimum
we use. Every presentation below is a finite structure, so the finite-carrier arguments of the
framework (its regime \textbf{A\_FIN}) apply to each presentation separately. The collection of
all presentations is infinite, and we never treat it as a finite carrier.

\subsection{Packages, lenses, audits}

\begin{convention}[Theory package]\label{conv:package}
A \emph{theory package} is a tuple $(Z,\,f,\,E,\,\mathcal A)$ consisting of a carrier $Z$; a
\emph{lens} $f\colon Z\to X$ recording what an observer sees of a carrier object; a
\emph{packaging} endomap $E\colon Z\to Z$, typically idempotent, whose fixed points are the
\emph{objects} the package recognizes; and an \emph{audit} $\mathcal A$, a certificate that
strictly decreases along the structure being audited. A predicate on $Z$ is \emph{visible} to
the lens if it is constant on every fiber of $f$.
\end{convention}

\subsection{The set-theoretic instrument}

\begin{definition}[Finite membership presentations]\label{def:presentation}
A \emph{finite pointed membership presentation} is a triple $P=(V,\to,v)$ where $V$ is a finite
set of nodes, ${\to}\subseteq V\times V$ is acyclic, and $v\in V$ is a distinguished
\emph{point}. We read $u\to w$ as ``$w$ is a member-candidate of $u$'' and write
$\Ch_P(u):=\{w: u\to w\}$. To keep the carrier a set, node sets are finite subsets of
$\mathbb N$. The carrier $Z$ is the set of all such presentations; it is countably infinite.
\end{definition}

In a finite graph, acyclicity is the same as well-foundedness: an infinite descending path in a
finite graph must repeat a node and therefore contains a cycle. This identification is the seed
of Foundation, and we return to it in \S\ref{ssec:foundation}.

\begin{definition}[The membership lens]\label{def:lens}
For $P=(V,\to,v)$ define the \emph{decoration} $\dec_P\colon V\to\HF$ by well-founded
recursion,
\[
  \dec_P(u)\;:=\;\{\dec_P(w): w\in\Ch_P(u)\},
\]
and set $f\colon Z\to\HF$, $f(P):=\dec_P(v)$. We call $f$ the \emph{membership-profile
lens}.\footnote{The term \emph{decoration} follows the graph-decoration viewpoint on membership
structures of \citet{Aczel1988}; on well-founded graphs it is the map collapsed by the Mostowski
lemma.}
\end{definition}

\begin{lemma}[Decoration is well posed; \inhouse]\label{lem:L1}
For every $P\in Z$ the decoration $\dec_P$ is uniquely defined, and $f(P)\in\HF$.
\end{lemma}

\begin{proof}
Finiteness and acyclicity give each node a finite height $h(u)$, the length of the longest
path starting at $u$, with $h(w)<h(u)$ whenever $u\to w$. Define $\dec_P$ by induction on
height: nodes of height $0$ are childless and decorate to $\varnothing$, and a node of height $n$
decorates to the finite set of the decorations of its children, all of lower height. The same
induction shows that each decoration is hereditarily finite and that the decoration is unique.
\end{proof}

\begin{definition}[Canonical presentation and packaging map]\label{def:E}
For $a\in\HF$ let $\TC(\{a\})$ be the least transitive set containing $a$ as an element, and let
$C(a):=(\TC(\{a\}),\in,a)$ be the \emph{canonical presentation} of $a$: its nodes are the
elements of $\TC(\{a\})$, named by their Ackermann codes in $\mathbb N$~\citep{Ackermann1937},
its edges are membership, and its point is $a$. The \emph{packaging map} is
\[
  E\colon Z\to Z,\qquad E(P):=C(f(P)).
\]
\end{definition}

In words, $E$ reads the membership profile of a presentation, forgets the presentation's
accidental shape, and presents the profile again in canonical form. Figure~\ref{fig:packaging}
shows a small example.

\begin{figure}[tbp]
\centering
\begin{tikzpicture}[
    font=\small,
    nd/.style={circle, draw, fill=white, minimum size=5.5mm, inner sep=0pt},
    pt/.style={circle, draw, thick, fill=blue!12, minimum size=5.5mm, inner sep=0pt},
    dl/.style={font=\scriptsize, text=blue!60!black},
    e/.style={-{Stealth[length=1.8mm]}, thick},
    >={Stealth[length=2mm]}]
  % --- left: an arbitrary presentation P ---
  \begin{scope}
    \node[pt] (r) at (0,2.4) {$r$};
    \node[nd] (u1) at (-1.3,1.2) {$u_1$};
    \node[nd] (u2) at (-0.35,1.2) {$u_2$};
    \node[nd] (w) at (1.1,1.2) {$w$};
    \node[nd] (u3) at (1.1,0) {$u_3$};
    \draw[e] (r) -- (u1); \draw[e] (r) -- (u2); \draw[e] (r) -- (w); \draw[e] (w) -- (u3);
    \node[dl, left=1pt of u1] {$\varnothing$};
    \node[dl, below=1pt of u2] {$\varnothing$};
    \node[dl, right=1pt of w] {$\{\varnothing\}$};
    \node[dl, right=1pt of u3] {$\varnothing$};
    \node[dl, right=2pt of r] {$\{\varnothing,\{\varnothing\}\}$};
    \node[font=\footnotesize] at (0,-0.8) {presentation $P$ (5 nodes)};
  \end{scope}
  % --- arrow ---
  \draw[->, very thick, gray] (2.9,1.3) -- node[above, font=\footnotesize, text=black]
        {$E=C\circ f$} node[below, font=\scriptsize, text=black, align=center]
        {read the profile,\\re-present it} (5.0,1.3);
  % --- right: canonical presentation C(2) ---
  \begin{scope}[xshift=7.2cm]
    \node[pt] (two) at (0,2.4) {$2$};
    \node[nd] (one) at (0.9,1.2) {$1$};
    \node[nd] (zero) at (0,0) {$0$};
    \draw[e] (two) -- (one); \draw[e] (two) -- (zero); \draw[e] (one) -- (zero);
    \node[font=\footnotesize, align=center] at (0,-0.8)
         {canonical $C(2)=E(P)$ (3 nodes)};
    \draw[->, gray, thick] (-0.25,2.75) arc (200:-20:0.32);
    \node[font=\scriptsize, anchor=south] at (0.05,3.05) {$E(E(P))=E(P)$};
  \end{scope}
\end{tikzpicture}
\caption{The packaging map on a small example. The point $r$ of $P$ (shaded) decorates to
$2=\{\varnothing,\{\varnothing\}\}$; the blue labels are the decorations $\dec_P$. The lens $f$ keeps only
this profile, and $E$ re-presents it as the canonical graph $C(2)$ on the transitive closure
$\TC(\{2\})=\{0,1,2\}$. The three childless nodes $u_1,u_2,u_3$ merge into the single node $0$.
Applying $E$ again changes nothing (Theorem~\ref{thm:p5}).}
\label{fig:packaging}
\end{figure}


\begin{definition}[The rank audit]\label{def:rank}
The audit is $\rho(P):=\rank(f(P))$, and for a node $u$ of $P$, $\rho_P(u):=\rank(\dec_P(u))$,
where $\rank(a)=\sup\{\rank(b)+1: b\in a\}$.
\end{definition}

\subsection{Extensionality as a choice of lens}

The membership lens identifies two presentations exactly when they have the same profile. This
identification is not a theorem about presentations; it is the choice of instrument.

\begin{lemma}[Lens equivalence; \inhouse]\label{lem:L2}
Define $P\sim Q$ iff $f(P)=f(Q)$. Then $\sim$ is an equivalence relation; each class contains the
canonical presentation of its profile, since $f(C(a))=a$; and on rooted extensional
presentations\footnote{Rooted: every node is reachable from the point. Extensional: distinct
nodes have distinct sets of children.} $P\sim Q$ holds exactly when $P$ and $Q$ are isomorphic as
pointed graphs.
\end{lemma}

\begin{proof}
The first clause is immediate. For $f(C(a))=a$, show $\dec_{C(a)}(s)=s$ for every
$s\in\TC(\{a\})$ by induction on rank. For the third clause, let $P$ be rooted and extensional.
The decoration $\dec_P$ is injective: if $\dec_P(u)=\dec_P(u')$, induct on $h(u)+h(u')$; equal
decorations give a matching between the children of $u$ and of $u'$ with equal decorations, the
inductive hypothesis identifies matched children, and extensionality then identifies $u$ and
$u'$. Rootedness makes the image of $\dec_P$ exactly $\TC(\{f(P)\})$, so $\dec_P$ is an
isomorphism of pointed graphs from $P$ onto $C(f(P))$. Two such presentations with the same
profile are therefore isomorphic to the same canonical presentation, and the converse is clear.
\end{proof}

\begin{remark}[What the lens forgets]\label{rem:lens-forgets}
Lemma~\ref{lem:L2} says that extensional identification is a genuine quotient: node names, the
multiplicity of routes, and graph shape beyond the profile are invisible to this instrument. In
Figure~\ref{fig:packaging} the three childless nodes are merged. Extensionality is therefore not
forced by the presentations themselves; it is the decision to look only through $f$. A different
instrument may lawfully keep apart what $f$ merges (countermodel \textbf{CM-1},
\S\ref{sec:atlas}).
\end{remark}

\subsection{The audit decreases along membership}

\begin{lemma}[Audit monotonicity; \inhouse]\label{lem:L3}
If $u\to w$ in $P$, then $\rho_P(w)<\rho_P(u)$. Equivalently, for $a\in\HF$ and $b\in a$,
$\rank(b)<\rank(a)$.
\end{lemma}

\begin{proof}
$\dec_P(w)\in\dec_P(u)$ by definition, and $\rank(a)=\sup\{\rank(b)+1: b\in a\}$.
\end{proof}

\paragraph{Role assignment.} In the framework's numbering of its six roles, the packaging map
$E$ is the packaging role (P5), whose fixed points are analysed in Theorem~\ref{thm:p5}; the rank
$\rho$ is the audit role (P6), whose exactness is Theorem~\ref{thm:tfae}; taking subsets of a
set (gating) is the constraint role (P2); forming images along functions is the descent or
transport role (P1); and the staging of sets by rank is the staging role (P4). The sixth role,
the route-mismatch diagnostic (P3), is not used to prove anything here.
